\begin{textsample}{NEG Dim 6 – global\_south\_2024\_08 – Score -1.00 – global\_south\_2024\_08/111601446.txt}  \label{ex:f6_neg_134}
Athletes from Israel and Palestine spar over the Gaza War in Paris PARIS : The war raging in their homeland was on the minds of Israeli and Palestinian judoists as they practiced on the tatami mats at the Olympic Games in Paris on Tuesday.
Two athletes competed against different opponents, one from each side, and lost.
However, after losing, they released contrasting statements emphasizing how intimate the competition was for the athletes and how challenging it had been for the organizers to establish an Olympic ceasefire following ten months of hostilities between Israeli forces and Palestinian Hamas militants in Gaza.
After losing his opening-round match in the under 81 kg division to Tajik Somon Makhmadbekov, Feras Badawi, one of eight Palestinian athletes competing at the Games, told reporters,"I think here at the Olympics, we are here to make peace, but if you are making war in our country and want to make peace here, it ' s like you have two faces."In light of the continuing crisis in Gaza, where at least 39,400 Palestinians @ @ @ @ @ @ @ @ @ @ retaliation for the 1,200-person attack on southern Israel on October 7, he stated he could never compete against Israeli athletes or shake their hands.
Shortly afterward, Israeli judoka Gili Sharir, competing in the women ' s under 63 kg division, was also thinking about losing to the current French Olympic champion, Clarisse Agbegnenou.
However, similar to Badawi, Gili Sharir was confronted with concerns regarding the wider political backdrop."We can take advantage of what is happening ; we can not ignore it.
One of the 88 Israeli athletes competing in Paris, Sharir said,"It was an honor for me to fight with the Israeli flag on my chest.
I believe in Israel, I love Israel."She expressed her wish that sports may take precedence over politics, stating that as an Israeli athlete, she was accustomed to people not shaking her hand.
The clashes at the Champs de Mars arena were the most recent illustration of how the war has thrown a shadow behind the scenes @ @ @ @ @ @ @ @ @ @ an Algerian judoka was disqualified from a possible match-up with an Israeli on Monday.
Prior to the Games, there were rumors that large-scale protests might occur in the periphery.
There was worry that notable acts expressing anti-Israeli sentiment may occur at competitions featuring Israeli athletes, although those have only occasionally occurred.
There have been flag waves for Palestine.
Some banners have demanded a"Free Palestine"or charged Israel with genocide, an accusation that Israel denies. 52 years after 11 Israelis were slain at the Munich Games in 1972, there have been some hateful chants at football games, and French authorities have launched an investigation into death threats against Israeli sportsmen.
The officials from both delegations have largely engaged in verbal sparring as they try to protect their positions.
Israeli officials chastised Palestinian competitors for wearing political emblems on their jackets at the opening ceremony, claiming that their actions violated the Olympic code.
The Palestinian Olympic Committee ( POC ) requested that Israel be suspended due to the war @ @ @ @ @ @ @ @ @ @ ) and the International Football Federation.
The POC said that this suspension was equivalent to Russia ' s suspension due to its invasion of Ukraine.
POC chief Jibril Rajoub, who served 17 years in Israeli prisons, detailed the severe consequences the war had had on Palestinian athletes on Tuesday and criticized what he called the IOC ' s"double standards"for breaching its own charter and failing to reply to his messages.
He specifically attacked Israel ' s judoka flagbearer Peter Paltchik over a media post and accused several Israeli athletes of glorifying the war on social media."Is he eligible to fly a flag at such a large, peaceful, humanitarian event?"Rajoub stated to the press."They ought to show Israel the red card."Rajoub ' s charges were refuted by the Israeli embassy in Paris, which described them as a"smear campaign"intended to"discredit and hurt him using lies and deception."Officials from the Israeli Olympic Committee did not hold back.
Rajoub was @ @ @ @ @ @ @ @ @ @"It ' s a ridiculous disgrace to use the arena for political or even more so for attacks of the Israeli athletes, but you know, people can choose who they want to cheer for,"she remarked.
However, we decide to focus on the positive.
The largest and most powerful symbol of Israel is the Israeli flag, which is what we have chosen to view.
When we observe it, we consider winning after October 7,"stated the former Olympian in judo.
PKKH ( Pakistan Ka Khuda Hafiz ) is a respected Pakistani news website offering accurate and comprehensive coverage of politics, culture, and global affairs.
With a commitment to objective reporting and insightful analysis, PKKH provides readers with a \textbf{well-rounded} perspective on current events.
The platform encourages community engagement and leverages digital technology to reach a wide audience, contributing to informed discussions and awareness.

% matched lemmas: well-rounded
\end{textsample}
